\documentclass[14pt,a4paper,titlepage]{report}
\usepackage[utf8]{inputenc}     % accents dans la source
\usepackage[T1]{fontenc}        % nouvelle norme latex
\usepackage[normalem]{ulem}     % 
\usepackage[french]{babel}      % supporter le français
\usepackage{verbatim}           % coder en VERBATIM
\usepackage{moreverb}                   % coder en VERBATIM
\usepackage[pdftex]{graphicx, color}
\graphicspath{{figs/}}
\DeclareGraphicsExtensions{.jpg,.png}
\usepackage{url}                %
\usepackage{float}
\usepackage{hyperref}           % utiliser les liens hypertextes
\usepackage{color}              % utiliser les couleur 
\usepackage{pifont}             % des symboles et des figures
\usepackage{fancybox}           % carrés 
\usepackage[a4paper,
  rmargin=2cm,
  lmargin=2cm,
  tmargin=2cm,
  bmargin=2cm]{geometry}  % mise en page 
\usepackage{listings}           % du code dans latex
\usepackage{fancyhdr}
\usepackage{bibunits}



\begin{document}
\begin{titlepage}
  \begin{flushleft}
    \vspace{3mm}
    \includegraphics[width=3cm]{upmclogo_m.png}\\
  \end{flushleft}
  \vfil
  \vfil
  \begin{center}
    \Huge\textbf{{PPAR}}\\
    \hrulefill\\
    \Huge\textit{{Probleme des N-CORPS}}\\
  \end{center}
  \hrulefill\\
  \vfil
  \noindent \textbf{ALAHYANE} \bsc{Rachid}
  \hfill 
  \href{mailto:rachid.alahyane@etu.upmc.fr}{rachid.alahyane@etu.upmc.fr}\\ 
  \textbf{CASTELLANOS} \bsc{Cuauhtemoc}
  \hfill 
  \href{mailto:cuauhtemoc.castellanos@etu.upmc.fr}{cuauhtemoc.castellanos@etu.upmc.fr}\\
  \vfil
  \begin{center}	
    \shadowbox{\footnotesize{\emph{PPAR 2009 - \today}}}
  \end{center}
\end{titlepage}





\lhead{PPAR 2009}
\rhead{Probleme des N-corps}
\renewcommand\headrulewidth{2pt}

\newpage
\pagestyle{empty}
\tableofcontents

\newpage
\chapter{Contexte et objectifs}
\thispagestyle{fancy}
\section{Contexte et objectifs}
\par Dans le cadre du module de programmation parallèle (PPAR)
du Master 2 systèmes et applications répartis (SAR), nous devons réaliser
la parallélisation d'un programme permettant de modéliser la dynamique des galaxies.
Ce problème d'astrophysique appelé aussi le problème des N-corps consiste à 
simuler les interactions dans un espace en 3 dimensions de N corps les uns avec
les autres.\\
\par Dans un premier temps, nous allons paralléliser le code séquentiel de manière 
à le rendre exécutable sur un petit cluster de machines de l'ARI. Pour cela nous 
utiliserons la bibliothèque MPI qui permet facilement de répartir les calculs sur 
plusieurs machines. Nous définirons le stockage en mémoire, l'architecture des
communications entre chaque machines ; nous recouvrirons également les temps de 
communication par le calcul et nous rendrons les communications persistantes.\\
\par Une fois cela effectué, nous utiliserons le fait que chaque noeud (machine)
possède un processeur multi-coeur, nous mettrons en place une parellélisation hybride
multi-processus/multi-threads en utilisant l'API OpenMP.\\
\par En fin, une étude comparative des performances sera fournie,
développant et expliquant les différents résultats des performances obtenues.\\



\chapter{Etat de l'art}
\thispagestyle{fancy}
%\section{Contexte et objectifs}
\par Dans le cadre du module de programmation parallèle (PPAR)
du Master 2 systèmes et applications répartis (SAR), nous devons réaliser
la parallélisation d'un programme permettant de modéliser la dynamique des galaxies.
Ce problème d'astrophysique appelé aussi le problème des N-corps consiste à 
simuler les interactions dans un espace en 3 dimensions de N corps les uns avec
les autres.\\
\par Dans un premier temps, nous allons paralléliser le code séquentiel de manière 
à le rendre exécutable sur un petit cluster de machines de l'ARI. Pour cela nous 
utiliserons la bibliothèque MPI qui permet facilement de répartir les calculs sur 
plusieurs machines. Nous définirons le stockage en mémoire, l'architecture des
communications entre chaque machines ; nous recouvrirons également les temps de 
communication par le calcul et nous rendrons les communications persistantes.\\
\par Une fois cela effectué, nous utiliserons le fait que chaque noeud (machine)
possède un processeur multi-coeur, nous mettrons en place une parellélisation hybride
multi-processus/multi-threads en utilisant l'API OpenMP.\\
\par En fin, une étude comparative des performances sera fournie,
développant et expliquant les différents résultats des performances obtenues.\\



\chapter{Architectures}
\thispagestyle{fancy}
%\section{Contexte et objectifs}
\par Dans le cadre du module de programmation parallèle (PPAR)
du Master 2 systèmes et applications répartis (SAR), nous devons réaliser
la parallélisation d'un programme permettant de modéliser la dynamique des galaxies.
Ce problème d'astrophysique appelé aussi le problème des N-corps consiste à 
simuler les interactions dans un espace en 3 dimensions de N corps les uns avec
les autres.\\
\par Dans un premier temps, nous allons paralléliser le code séquentiel de manière 
à le rendre exécutable sur un petit cluster de machines de l'ARI. Pour cela nous 
utiliserons la bibliothèque MPI qui permet facilement de répartir les calculs sur 
plusieurs machines. Nous définirons le stockage en mémoire, l'architecture des
communications entre chaque machines ; nous recouvrirons également les temps de 
communication par le calcul et nous rendrons les communications persistantes.\\
\par Une fois cela effectué, nous utiliserons le fait que chaque noeud (machine)
possède un processeur multi-coeur, nous mettrons en place une parellélisation hybride
multi-processus/multi-threads en utilisant l'API OpenMP.\\
\par En fin, une étude comparative des performances sera fournie,
développant et expliquant les différents résultats des performances obtenues.\\



\chapter{Statistiques}
\thispagestyle{fancy}
%\section{Contexte et objectifs}
\par Dans le cadre du module de programmation parallèle (PPAR)
du Master 2 systèmes et applications répartis (SAR), nous devons réaliser
la parallélisation d'un programme permettant de modéliser la dynamique des galaxies.
Ce problème d'astrophysique appelé aussi le problème des N-corps consiste à 
simuler les interactions dans un espace en 3 dimensions de N corps les uns avec
les autres.\\
\par Dans un premier temps, nous allons paralléliser le code séquentiel de manière 
à le rendre exécutable sur un petit cluster de machines de l'ARI. Pour cela nous 
utiliserons la bibliothèque MPI qui permet facilement de répartir les calculs sur 
plusieurs machines. Nous définirons le stockage en mémoire, l'architecture des
communications entre chaque machines ; nous recouvrirons également les temps de 
communication par le calcul et nous rendrons les communications persistantes.\\
\par Une fois cela effectué, nous utiliserons le fait que chaque noeud (machine)
possède un processeur multi-coeur, nous mettrons en place une parellélisation hybride
multi-processus/multi-threads en utilisant l'API OpenMP.\\
\par En fin, une étude comparative des performances sera fournie,
développant et expliquant les différents résultats des performances obtenues.\\



\chapter{Problèmes rencontrés }
\thispagestyle{fancy}
%\section{Contexte et objectifs}
\par Dans le cadre du module de programmation parallèle (PPAR)
du Master 2 systèmes et applications répartis (SAR), nous devons réaliser
la parallélisation d'un programme permettant de modéliser la dynamique des galaxies.
Ce problème d'astrophysique appelé aussi le problème des N-corps consiste à 
simuler les interactions dans un espace en 3 dimensions de N corps les uns avec
les autres.\\
\par Dans un premier temps, nous allons paralléliser le code séquentiel de manière 
à le rendre exécutable sur un petit cluster de machines de l'ARI. Pour cela nous 
utiliserons la bibliothèque MPI qui permet facilement de répartir les calculs sur 
plusieurs machines. Nous définirons le stockage en mémoire, l'architecture des
communications entre chaque machines ; nous recouvrirons également les temps de 
communication par le calcul et nous rendrons les communications persistantes.\\
\par Une fois cela effectué, nous utiliserons le fait que chaque noeud (machine)
possède un processeur multi-coeur, nous mettrons en place une parellélisation hybride
multi-processus/multi-threads en utilisant l'API OpenMP.\\
\par En fin, une étude comparative des performances sera fournie,
développant et expliquant les différents résultats des performances obtenues.\\


\chapter{Conclusion}
\thispagestyle{fancy}
%\section{Contexte et objectifs}
\par Dans le cadre du module de programmation parallèle (PPAR)
du Master 2 systèmes et applications répartis (SAR), nous devons réaliser
la parallélisation d'un programme permettant de modéliser la dynamique des galaxies.
Ce problème d'astrophysique appelé aussi le problème des N-corps consiste à 
simuler les interactions dans un espace en 3 dimensions de N corps les uns avec
les autres.\\
\par Dans un premier temps, nous allons paralléliser le code séquentiel de manière 
à le rendre exécutable sur un petit cluster de machines de l'ARI. Pour cela nous 
utiliserons la bibliothèque MPI qui permet facilement de répartir les calculs sur 
plusieurs machines. Nous définirons le stockage en mémoire, l'architecture des
communications entre chaque machines ; nous recouvrirons également les temps de 
communication par le calcul et nous rendrons les communications persistantes.\\
\par Une fois cela effectué, nous utiliserons le fait que chaque noeud (machine)
possède un processeur multi-coeur, nous mettrons en place une parellélisation hybride
multi-processus/multi-threads en utilisant l'API OpenMP.\\
\par En fin, une étude comparative des performances sera fournie,
développant et expliquant les différents résultats des performances obtenues.\\



\end{document}
